\section*{NeurIPS Paper Checklist}

\begin{enumerate}

\item {\bf Claims}
    \item[] Question: Do the main claims made in the abstract and introduction accurately reflect the paper's contributions and scope?
    \item[] Answer: \answerYes{}
    \item[] Justification: The abstract and Section~\ref{sec:introduction} state that refutation is sound relative to the declared pack, distinguish \textsc{Unknown} from refutation, and identify the mechanized and paper-only proof fragments.

\item {\bf Limitations}
    \item[] Question: Does the paper discuss the limitations of the work performed by the authors?
    \item[] Answer: \answerYes{}
    \item[] Justification: The Limitations and Future Work section discusses relative soundness, static topology, abstraction coarseness, mechanization gaps, \new{low recall on real files with LLM compaction, and the size of the held-out sets.}

\item {\bf Theory assumptions and proofs}
    \item[] Question: For each theoretical result, does the paper provide the full set of assumptions and a complete (and correct) proof?
    \item[] Answer: \answerYes{}
    \item[] Justification: Appendix~\ref{app:metatheory} states the simulation, stability, commutativity, finiteness, and topology assumptions with proofs or proof sketches; the artifact checks the explicitly identified Coq fragments.

\item {\bf Experimental result reproducibility}
    \item[] Question: Does the paper fully disclose all the information needed to reproduce the main experimental results of the paper to the extent that it affects the main claims and/or conclusions of the paper (regardless of whether the code and data are provided or not)?
    \item[] Answer: \answerYes{}
    \item[] Justification: Sections~\ref{sec:implementation}--\ref{sec:evaluation} describe the deterministic checker, verdict semantics, corpus construction, binary accounting, and separate treatment of abstentions.

\item {\bf Open access to data and code}
    \item[] Question: Does the paper provide open access to the data and code, with sufficient instructions to faithfully reproduce the main experimental results, as described in supplemental material?
    \item[] Answer: \answerYes{}
    \item[] Justification: The \new{public} artifact contains the checker, synthetic corpora, proof developments, pinned continuous-integration workflow, and reproduction commands\new{, plus the evaluation sets with their labels and the reports of the model runs.}

\item {\bf Experimental setting/details}
    \item[] Question: Does the paper specify all the training and test details (e.g., data splits, hyperparameters, how they were chosen, type of optimizer) necessary to understand the results?
    \item[] Answer: \new{\answerYes{}}
    \item[] Justification: \new{No model is trained. Section~\ref{sec:evaluation} names the model (gpt-5.4), the number of compaction samples per file, the sets (set2 for development, set3 held out and labelled blind) and the baselines, and reports where the results fall short.}

\item {\bf Experiment statistical significance}
    \item[] Question: Does the paper report error bars suitably and correctly defined or other appropriate information about the statistical significance of the experiments?
    \item[] Answer: \new{\answerYes{}}
    \item[] Justification: \new{The checker is deterministic, so the hand-written corpus has no sampling variance. The model-based results report counts with their denominators and 95\% confidence intervals (Figures~\ref{fig:usecases} and~\ref{fig:cost}), and agreement across samples.}

\item {\bf Experiments compute resources}
    \item[] Question: For each experiment, does the paper provide sufficient information on the computer resources (type of compute workers, memory, time of execution) needed to reproduce the experiments?
    \item[] Answer: \answerYes{}
    \item[] Justification: \new{The checker runs on a CPU (a median of 32\,ms per decision in Q2) with a 10-second solver timeout per query. Model runs use hosted gpt-5.4; their token counts are reported in Section~\ref{sec:evaluation} and Appendix~\ref{app:tokens}.}

\item {\bf Code of ethics}
    \item[] Question: Does the research conducted in the paper conform, in every respect, with the NeurIPS Code of Ethics \url{https://neurips.cc/public/EthicsGuidelines}?
    \item[] Answer: \answerYes{}
    \item[] Justification: The work uses synthetic packs \new{and public skill and agent files}, no personal data or human subjects, and makes bounded claims with explicit trust assumptions and limitations. \new{Benchmark files are read as text and never executed.}

\item {\bf Broader impacts}
    \item[] Question: Does the paper discuss both potential positive societal impacts and negative societal impacts of the work performed?
    \item[] Answer: \answerYes{}
    \item[] Justification: The Broader impact paragraph discusses reduced wasted execution as well as false confidence from inaccurate packs and the human/runtime mitigations required.

\item {\bf Safeguards}
    \item[] Question: Does the paper describe safeguards that have been put in place for responsible release of data or models that have a high risk for misuse (e.g., pre-trained language models, image generators, or scraped datasets)?
    \item[] Answer: \answerNA{}
    \item[] Justification: The artifact releases no trained model, scraped dataset, or other high-risk asset; it contains a deterministic checker and synthetic examples.

\item {\bf Licenses for existing assets}
    \item[] Question: Are the creators or original owners of assets (e.g., code, data, models), used in the paper, properly credited and are the license and terms of use explicitly mentioned and properly respected?
    \item[] Answer: \new{\answerYes{}}
    \item[] Justification: \new{The real files come from public repositories and are pinned by commit and sha256 in the artifact; they are used only as test inputs. The models are used through their providers' services. Prior methods and software foundations are credited through citations.}

\item {\bf New assets}
    \item[] Question: Are new assets introduced in the paper well documented and is the documentation provided alongside the assets?
    \item[] Answer: \answerYes{}
    \item[] Justification: The supplemental artifact documents the checker interface, pack schema, synthetic corpora, proof scope, limitations, and reproduction workflow.

\item {\bf Crowdsourcing and research with human subjects}
    \item[] Question: For crowdsourcing experiments and research with human subjects, does the paper include the full text of instructions given to participants and screenshots, if applicable, as well as details about compensation (if any)?
    \item[] Answer: \answerNA{}
    \item[] Justification: The research involves neither crowdsourcing nor human-subject experiments.

\item {\bf Institutional review board (IRB) approvals or equivalent for research with human subjects}
    \item[] Question: Does the paper describe potential risks incurred by study participants, whether such risks were disclosed to the subjects, and whether Institutional Review Board (IRB) approvals (or an equivalent approval/review based on the requirements of your country or institution) were obtained?
    \item[] Answer: \answerNA{}
    \item[] Justification: The research involves no study participants or human-subject data.

\item {\bf Declaration of LLM usage}
    \item[] Question: Does the paper describe the usage of LLMs if it is an important, original, or non-standard component of the core methods in this research? Note that if the LLM is used only for writing, editing, or formatting purposes and does not impact the core methodology, scientific rigor, or originality of the research, declaration is not required.
    \item[] Answer: \new{\answerYes{}}
    \item[] Justification: \new{LLMs are part of the evaluated pipeline: gpt-5.4 compacts files into packs and gives the direct-verdict baseline, DeepSeek-V4-Pro is a second model family, and a Claude model judges pack fidelity (Section~\ref{sec:evaluation}).} The paper's core method does not depend on an LLM. The main contributions---the achievability type discipline, its formal metatheory, and the deterministic checker---are defined and verified independently of LLM inference. The LLM is used only as a front-end to compact a natural-language skill into the formal pack consumed by the checker; the same formal pack can be provided directly without an LLM. The formal guarantees therefore do not depend on the behavior of any particular LLM\new{; the measured recall and false-refutation rates of the full pipeline do}.

\end{enumerate}
